\documentclass[11pt,a4paper]{article}

\usepackage[T1]{fontenc}
\usepackage[utf8]{inputenc}
\usepackage[french]{babel}
\usepackage{lmodern}
\usepackage{amsmath,amssymb}
\usepackage{geometry}
\usepackage{tikz}
\usepackage{pgfplots}
\usepackage{enumitem}
\usepackage{xcolor}
\usepackage{fancyhdr}
\usepackage[many]{tcolorbox}

\pgfplotsset{compat=1.18}
\geometry{top=1.75cm,bottom=1.15cm,left=1.25cm,right=1.25cm,headheight=14pt,headsep=0.35cm}
\setlength{\parindent}{0pt}
\setlength{\parskip}{2pt}

\definecolor{courberouge}{RGB}{200,55,55}
\definecolor{tangentebleue}{RGB}{45,105,180}
\definecolor{tangenteverte}{RGB}{55,125,60}
\definecolor{violetpoint}{RGB}{90,55,170}
\definecolor{orangepoint}{RGB}{220,120,35}
\definecolor{corrigecolor}{RGB}{30,110,70}

\pagestyle{fancy}
\fancyhf{}
\lhead{Contrôle}
\chead{Tangentes et automatismes}
\rhead{Première spécialité}
\cfoot{\thepage}

\newcommand{\choixquatre}[4]{%
\begin{tabular}{@{}p{0.23\linewidth}p{0.23\linewidth}p{0.23\linewidth}p{0.23\linewidth}@{}}
\textbf{A.} #1 & \textbf{B.} #2 & \textbf{C.} #3 & \textbf{D.} #4
\end{tabular}}


\begin{document}


\begin{center}
    {\Large \textbf{Contrôle - Tangentes et automatismes}}\par
    \vspace{0.15cm}
    \textbf{Nom :} \dotfill \hspace{1cm} \textbf{Prénom :} \dotfill
\end{center}

\section*{Exercice 1 - Tangentes et nombre dérivé}

\begin{tcolorbox}[colback=gray!6,colframe=black!45,arc=2mm,left=2mm,right=2mm,top=1mm,bottom=1mm]
On considère une fonction $f$ définie sur l'intervalle représenté ci-dessous. On note $\mathcal C_f$ sa courbe représentative dans un repère.

La droite $\mathcal T_{-1}$ est la tangente à $\mathcal C_f$ au point d'abscisse $-1$.
La droite $\mathcal T_{0{,}5}$ est la tangente à $\mathcal C_f$ au point d'abscisse $0{,}5$.

\textbf{Aucune expression de la fonction $f$ n'est donnée. Les réponses doivent être obtenues à partir du graphique.}
\end{tcolorbox}

\begin{center}
\begin{tikzpicture}
\begin{axis}[
    width=16.2cm,
    height=9.6cm,
    xmin=-2.2, xmax=2.2,
    ymin=-0.3, ymax=3.35,
    axis lines=middle,
    axis line style={black,very thick},
    xlabel={$x$}, ylabel={$y$},
    xlabel style={at={(axis description cs:1,0.09)},anchor=west},
    ylabel style={at={(axis description cs:0.505,1)},anchor=south},
    xtick={-2,-1.5,-1,-0.5,0,0.5,1,1.5,2},
    ytick={0,0.5,1,1.5,2,2.5,3},
    minor x tick num=1,
    minor y tick num=1,
    grid=both,
    major grid style={line width=0.45pt,draw=black!35},
    minor grid style={line width=0.25pt,draw=black!14},
    tick label style={font=\small},
    clip=false
]

% Courbe de f : l'expression n'est pas donnée aux élèves dans l'énoncé.
\draw[courberouge, very thick]
    (axis cs:-2,2.6)
    .. controls (axis cs:-1.7667,2.0867) and (axis cs:-1.5333,0.5) .. (axis cs:-1.3,0.5)
    .. controls (axis cs:-1.2,0.5) and (axis cs:-1.1,0.8) .. (axis cs:-1,1)
    .. controls (axis cs:-0.7333,1.5333) and (axis cs:-0.4667,2.5) .. (axis cs:-0.2,2.5)
    .. controls (axis cs:0.0333,2.5) and (axis cs:0.2667,1.7333) .. (axis cs:0.5,1.5)
    .. controls (axis cs:0.7667,1.2333) and (axis cs:1.0333,0.8) .. (axis cs:1.3,0.8)
    .. controls (axis cs:1.5333,0.8) and (axis cs:1.7667,1.9333) .. (axis cs:2,2.4);

% Tangente en x=-1 : y=2x+3
\addplot[domain=-1.65:0.15, samples=2, very thick, tangentebleue]
    {2*x+3};

% Tangente en x=0,5 : y=-x+2
\addplot[domain=-0.85:2.1, samples=2, very thick, tangenteverte]
    {-x+2};

% Points repères visibles sur les tangentes
\draw[violetpoint, line width=3pt] (axis cs:-1.06,1) -- (axis cs:-0.94,1) (axis cs:-1,0.94) -- (axis cs:-1,1.06);
\draw[black, line width=3pt] (axis cs:-0.06,3) -- (axis cs:0.06,3) (axis cs:0,2.94) -- (axis cs:0,3.06);
\draw[orangepoint, line width=3pt] (axis cs:0.44,1.5) -- (axis cs:0.56,1.5) (axis cs:0.5,1.44) -- (axis cs:0.5,1.56);
\draw[black, line width=3pt] (axis cs:0.94,1) -- (axis cs:1.06,1) (axis cs:1,0.94) -- (axis cs:1,1.06);

\end{axis}
\end{tikzpicture}
\end{center}

\begin{enumerate}[label=\textbf{\arabic*)}, leftmargin=0.75cm, itemsep=0.18cm, topsep=0.1cm]
    \item Sur le graphique, repérer $\mathcal C_f$, $\mathcal T_{-1}$ et $\mathcal T_{0{,}5}$. Écrire directement les trois noms au bon endroit sur le graphique.
    \item Calculer le coefficient directeur de la tangente $\mathcal T_{-1}$ à l'aide de deux points lisibles sur le graphique.
    \item Donner les valeurs de $f'(-1)$ et de $f(-1)$.
    \item Déterminer l'équation réduite complète de la tangente $\mathcal T_{-1}$.
    \item Déterminer la valeur de $f'(0{,}5)$.
    \item Déterminer l'équation réduite complète de la tangente $\mathcal T_{0{,}5}$.
\end{enumerate}

\newpage

\section*{Exercice 2 - Automatismes}

Pour chaque question, entourer la bonne réponse.

\begin{enumerate}[label=\textbf{\arabic*)}, leftmargin=0.75cm, itemsep=0.55cm, topsep=0.2cm]

\item Une durée de $2{,}2$ heures correspond à :\par\smallskip
\choixquatre{$132$ minutes}{$142$ minutes}{$220$ minutes}{$122$ minutes}

\item Le prix d'un sac a baissé de $28\,\%$. Il coûte maintenant $192$ euros.\par
Le prix d'avant, en euros, est donné par le calcul :\par\smallskip
\choixquatre{$192\times \dfrac{28}{72}$}{$192\times 0{,}72$}{$\dfrac{192}{1-0{,}28}$}{$\dfrac{192}{1+\dfrac{28}{100}}$}

\item $38^2-37^2$ est égal à :\par\smallskip
\choixquatre{$75$}{$-75$}{$-1$}{$1$}

\item Soit $n$ un entier. À quelle expression est égale $(4^n)^2$ ?\par\smallskip
\choixquatre{$4^{n^2}$}{Aucune de ces propositions}{$16^n$}{$8^n$}

\item On sait que :
\[
0{,}1\times 1{,}5=0{,}15,\qquad 1{,}9\times 0{,}5=0{,}95,
\]
\[
1{,}9\times 1{,}5=2{,}85,\qquad 0{,}9\times 0{,}5=0{,}45.
\]
En utilisant l'un des résultats précédents, déterminer le taux global d'évolution d'un article qui augmente de $90\,\%$ dans un premier temps, puis qui augmente de $50\,\%$ dans un second temps.\par\smallskip
\choixquatre{$185\,\%$}{$140\,\%$}{$95\,\%$}{$285\,\%$}

\item On considère les trois fonctions définies par :
\[
f_1:x\longmapsto x^2-(x+5)(x-3),\qquad
f_2:x\longmapsto 4\sqrt{x}+4,\qquad
f_3:x\longmapsto \frac{4}{x}-3.
\]
On peut affirmer que :\par\smallskip
\begin{tabular}{@{}p{0.96\linewidth}@{}}
\textbf{A.} Toutes ces fonctions sont affines.\\
\textbf{B.} Uniquement la fonction $f_1$ est affine.\\
\textbf{C.} Aucune de ces fonctions n'est affine.\\
\textbf{D.} Uniquement les fonctions $f_2$ et $f_3$ sont affines.
\end{tabular}

\item Un prix a doublé. Cela signifie que le prix a augmenté de :\par\smallskip
\choixquatre{$100\,\%$}{$200\,\%$}{$120\,\%$}{$20\,\%$}

\item Le double de $4^{40}$ est égal à :\par\smallskip
\choixquatre{$4^{80}$}{$8^{40}$}{$2^{81}$}{$2^{41}$}

\end{enumerate}



\clearpage
\fancyhf{}
\lhead{Corrigé}
\chead{Tangentes et automatismes}
\rhead{Première spécialité}
\cfoot{\thepage}
\setcounter{page}{1}


\begin{center}
    {\Large \textbf{Corrigé - Tangentes et automatismes}}
\end{center}

\section*{Exercice 1 - Tangentes et nombre dérivé}

\begin{center}
\begin{tikzpicture}
\begin{axis}[
    width=15.8cm,
    height=8.9cm,
    xmin=-2.2, xmax=2.2,
    ymin=-0.3, ymax=3.35,
    axis lines=middle,
    axis line style={black,very thick},
    xlabel={$x$}, ylabel={$y$},
    xlabel style={at={(axis description cs:1,0.09)},anchor=west},
    ylabel style={at={(axis description cs:0.505,1)},anchor=south},
    xtick={-2,-1.5,-1,-0.5,0,0.5,1,1.5,2},
    ytick={0,0.5,1,1.5,2,2.5,3},
    minor x tick num=1,
    minor y tick num=1,
    grid=both,
    major grid style={line width=0.45pt,draw=black!35},
    minor grid style={line width=0.25pt,draw=black!14},
    tick label style={font=\small},
    clip=false
]

% Courbe de f
\draw[courberouge, very thick]
    (axis cs:-2,2.6)
    .. controls (axis cs:-1.7667,2.0867) and (axis cs:-1.5333,0.5) .. (axis cs:-1.3,0.5)
    .. controls (axis cs:-1.2,0.5) and (axis cs:-1.1,0.8) .. (axis cs:-1,1)
    .. controls (axis cs:-0.7333,1.5333) and (axis cs:-0.4667,2.5) .. (axis cs:-0.2,2.5)
    .. controls (axis cs:0.0333,2.5) and (axis cs:0.2667,1.7333) .. (axis cs:0.5,1.5)
    .. controls (axis cs:0.7667,1.2333) and (axis cs:1.0333,0.8) .. (axis cs:1.3,0.8)
    .. controls (axis cs:1.5333,0.8) and (axis cs:1.7667,1.9333) .. (axis cs:2,2.4);

% Tangente en x=-1 : y=2x+3
\addplot[domain=-1.65:0.15, samples=2, very thick, tangentebleue]
    {2*x+3};

% Tangente en x=0,5 : y=-x+2
\addplot[domain=-0.85:2.1, samples=2, very thick, tangenteverte]
    {-x+2};

% Points repères visibles sur les tangentes
\draw[violetpoint, line width=3pt] (axis cs:-1.06,1) -- (axis cs:-0.94,1) (axis cs:-1,0.94) -- (axis cs:-1,1.06);
\draw[black, line width=3pt] (axis cs:-0.06,3) -- (axis cs:0.06,3) (axis cs:0,2.94) -- (axis cs:0,3.06);
\draw[orangepoint, line width=3pt] (axis cs:0.44,1.5) -- (axis cs:0.56,1.5) (axis cs:0.5,1.44) -- (axis cs:0.5,1.56);
\draw[black, line width=3pt] (axis cs:0.94,1) -- (axis cs:1.06,1) (axis cs:1,0.94) -- (axis cs:1,1.06);

\node[courberouge, font=\bfseries] at (axis cs:-1.85,2.95) {$\mathcal C_f$};
\node[tangentebleue, font=\bfseries] at (axis cs:-1.55,0.25) {$\mathcal T_{-1}$};
\node[tangenteverte, font=\bfseries] at (axis cs:1.65,0.55) {$\mathcal T_{0{,}5}$};

\end{axis}
\end{tikzpicture}
\end{center}

\begin{enumerate}[label=\textbf{\arabic*)}, leftmargin=0.75cm, itemsep=0.35cm]
    \item On repère : la courbe rouge est $\mathcal C_f$, la droite bleue est $\mathcal T_{-1}$, et la droite verte est $\mathcal T_{0{,}5}$.

    \item Pour calculer le coefficient directeur de $\mathcal T_{-1}$, on peut choisir deux points lisibles sur la droite bleue :
    \[
    A(-1;1) \qquad \text{et} \qquad B(0;3).
    \]
    Le coefficient directeur de la droite $(AB)$ est :
    \[
    m=\frac{\operatorname{ord}_B-\operatorname{ord}_A}{\operatorname{abs}_B-\operatorname{abs}_A}.
    \]
    Donc :
    \[
    m=\frac{3-1}{0-(-1)}=\frac{2}{1}=2.
    \]
    Ainsi, le coefficient directeur de $\mathcal T_{-1}$ est $2$.

    \item $f'(-1)$ est le coefficient directeur de la tangente à $\mathcal C_f$ au point d'abscisse $-1$.
    D'après la question précédente :
    \[
    \boxed{f'(-1)=2}.
    \]
    De plus, le point de $\mathcal C_f$ d'abscisse $-1$ a pour ordonnée $1$, donc :
    \[
    \boxed{f(-1)=1}.
    \]

    \item La tangente $\mathcal T_{-1}$ passe par le point $A(-1;1)$ et a pour coefficient directeur $2$.
    Son équation est donc :
    \[
    y=2(x+1)+1.
    \]
    On développe :
    \[
    y=2x+2+1.
    \]
    Ainsi :
    \[
    \boxed{\mathcal T_{-1}:\ y=2x+3}.
    \]

    \item $f'(0{,}5)$ est le coefficient directeur de la tangente à $\mathcal C_f$ au point d'abscisse $0{,}5$.

    Pour calculer ce coefficient directeur, on peut choisir deux points lisibles sur la droite verte :
    \[
    C(0{,}5;1{,}5) \qquad \text{et} \qquad D(1;1).
    \]
    On obtient :
    \[
    m=\frac{1-1{,}5}{1-0{,}5}=\frac{-0{,}5}{0{,}5}=-1.
    \]
    Donc :
    \[
    \boxed{f'(0{,}5)=-1}.
    \]

    \item On lit aussi sur le graphique :
    \[
    f(0{,}5)=1{,}5.
    \]
    La tangente $\mathcal T_{0{,}5}$ passe donc par le point $C(0{,}5;1{,}5)$ et a pour coefficient directeur $-1$.
    Son équation est :
    \[
    y=-1(x-0{,}5)+1{,}5.
    \]
    Donc :
    \[
    y=-x+0{,}5+1{,}5.
    \]
    Ainsi :
    \[
    \boxed{\mathcal T_{0{,}5}:\ y=-x+2}.
    \]
\end{enumerate}

\newpage

\section*{Exercice 2 - Automatismes}

\begin{enumerate}[label=\textbf{\arabic*)}, leftmargin=0.75cm, itemsep=0.55cm]

\item Une durée de $2{,}2$ heures correspond à :
\[
2{,}2\times 60=132.
\]
\[
\boxed{132\text{ minutes}} \qquad \text{Réponse A.}
\]

\item Le prix a baissé de $28\,\%$, donc il a été multiplié par :
\[
1-0{,}28=0{,}72.
\]
Si l'ancien prix est $P$, alors :
\[
0{,}72P=192.
\]
Donc :
\[
P=\frac{192}{0{,}72}=\frac{192}{1-0{,}28}.
\]
\[
\boxed{\dfrac{192}{1-0{,}28}} \qquad \text{Réponse C.}
\]

\item On utilise l'identité remarquable $a^2-b^2=(a-b)(a+b)$ :
\[
38^2-37^2=(38-37)(38+37)=1\times75=75.
\]
\[
\boxed{75} \qquad \text{Réponse A.}
\]

\item On utilise la règle $(a^n)^m=a^{nm}$ :
\[
(4^n)^2=4^{2n}=(4^2)^n=16^n.
\]
\[
\boxed{16^n} \qquad \text{Réponse C.}
\]

\item Une augmentation de $90\,\%$ correspond à un coefficient multiplicateur de $1{,}9$.
Une augmentation de $50\,\%$ correspond à un coefficient multiplicateur de $1{,}5$.

Le coefficient multiplicateur global est donc :
\[
1{,}9\times1{,}5=2{,}85.
\]
Cela correspond à une augmentation de :
\[
2{,}85-1=1{,}85,
\]
soit :
\[
\boxed{185\,\%} \qquad \text{Réponse A.}
\]

\item On simplifie $f_1$ :
\[
f_1(x)=x^2-(x+5)(x-3).
\]
Or :
\[
(x+5)(x-3)=x^2+2x-15.
\]
Donc :
\[
f_1(x)=x^2-(x^2+2x-15)=-2x+15.
\]
Ainsi $f_1$ est affine.

En revanche, $f_2:x\mapsto4\sqrt{x}+4$ et $f_3:x\mapsto\dfrac4x-3$ ne sont pas affines.

\[
\boxed{\text{Uniquement la fonction }f_1\text{ est affine.}} \qquad \text{Réponse B.}
\]

\item Si un prix double, il est multiplié par $2$.
Le taux d'évolution correspondant est :
\[
2-1=1,
\]
soit :
\[
\boxed{100\,\%} \qquad \text{Réponse A.}
\]

\item On écrit $4^{40}$ avec des puissances de $2$ :
\[
4^{40}=(2^2)^{40}=2^{80}.
\]
Donc :
\[
2\times4^{40}=2\times2^{80}=2^{81}.
\]
\[
\boxed{2^{81}} \qquad \text{Réponse C.}
\]

\end{enumerate}


\end{document}
